%% kapa: exam=1 round=19 year=2008 period=2 form=A subject=economics subject_ko=경제원론 seq=33 number=33
%% source: Q-Net 감정평가사 기출문제 (cst00309) artlSeq=1001962 file=extracted/1008/1001962/2008년 제19회 감정평가사 1차시험 기출문제/2교시.hwp
%% conversion: kordoc 4.24.51 -> tools/gen1.py (auto); review=proofread-sonnet; flags=none
\begin{question}{33}
영희는 하루에 아이스크림을 100개 팔고 있다. 영희의 하루 총비용은 100,000원이고, 이 중 고정비용은 10,000원이다. 완전경쟁시장에서 이윤극대화를 추구하는 영희가 단기에는 이 아이스크림 가게를 운영하지만 장기에는 폐업해야 하는 아이스크림 1개당 가격(P)의 구간으로 다음 중 가장 적절한 것은?
\begin{choices}
  \item 600원 $<$ P $<$ 700원
  \item 700원 $<$ P $<$ 800원
  \item 800원 $<$ P $<$ 900원
  \item 900원 $<$ P $<$ 1,000원
  \item 1,000원 $<$ P $<$ 1,100원
\end{choices}
\end{question}
\answer{4}
